% ============================================================
% Formulario, argomento 26: codifica di sorgente (dispensa, capitolo 26)
% ============================================================
\argomento{La codifica di sorgente}{Dispensa: capitolo 26}

\voce{codice}{Codice di sorgente e lunghezza media}{sorgente senza memoria (DMS): lettere i.i.d. con pmf $p_X$}
  {C : x \in \mathcal{X} \;\longmapsto\; C(x) \in \{0, 1\}^{l(x)} \qquad L = E[l(X)] = \sum_{x} p_X(x)\, l(x)}
  {\variante{sulle sequenze si concatena}{C(x_1 x_2 \cdots x_n) = C(x_1)\, C(x_2) \cdots C(x_n)}}

\voce{prefisso}{Codici a prefisso (istantanei)}{nessuna parola codice è l'inizio di un'altra}
  {C(x) \text{ non è prefisso di } C(x') \text{ per ogni } x \neq x'}
  {\variante{univocamente decifrabile: ogni stringa di bit si legge in un solo modo}{}
   \variante{a prefisso $\Rightarrow$ univocamente decifrabile: si decodifica leggendo da sinistra}{}
   \variante{albero binario: 0 a sinistra, 1 a destra, le parole sulle foglie}{}}

\voce{kraft}{Disuguaglianza di Kraft}{esiste un codice a prefisso con lunghezze $l(x)$ se e solo se}
  {\sum_{x \in \mathcal{X}} 2^{-l(x)} \leq 1}
  {\variante{vale anche per i codici univocamente decifrabili (McMillan)}{}
   \variante{con l'uguaglianza l'albero è pieno}{}}

\voce{limiteentropico}{Limite entropico}{per ogni codice univocamente decifrabile}
  {L = \sum_{x} p_X(x)\, l(x) \geq H(X)}
  {\variante{uguaglianza se e solo se le probabilità sono diadiche}{p_X(x) = 2^{-l(x)} \quad \text{per ogni } x}}

\voce{shannoncodice}{Codice di Shannon}{una parola lunga quanto l'informazione, arrotondata per eccesso}
  {l(x_i) = \left\lceil \log_2 \frac{1}{p_X(x_i)} \right\rceil \qquad H(X) \leq L < H(X) + 1}
  {}

\voce{huffman}{Algoritmo di Huffman}{il codice a prefisso con la minima lunghezza media}
  {\text{si uniscono ripetutamente i due simboli meno probabili}}
  {\variante{1) si ordinano i simboli per probabilità decrescente}{}
   \variante{2) i due meno probabili diventano un macrosimbolo, con la somma delle probabilità}{}
   \variante{3) si ripete fino alla radice, che ha probabilità 1}{}
   \variante{4) 0 e 1 ai due rami di ogni unione: la parola è il cammino dalla radice alla foglia}{}
   \variante{non è unico, ma la lunghezza media sì; le due lettere meno probabili differiscono solo nell'ultimo
     bit}{}
   \variante{se il limite entropico si può raggiungere, Huffman lo raggiunge}{}}

\voce{blocchi}{Codifica a blocchi}{codice di Shannon su blocchi di $k$ lettere di una sorgente senza memoria \oltre}
  {H(X) \leq \frac{L_k}{k} < H(X) + \frac{1}{k}}
  {\variante{$L_k$ è la lunghezza media per blocco; l'entropia di un blocco è}{k\, H(X)}}
